\section{Introduction}\label{intro:section}
For a finite set $U\subset\R^2$, let
\[
 u(U)=\#\{\{x,y\}\subset U:|x-y|=1\},\qquad
 u(n)=\max_{|U|=n}u(U),
\]
where the pairs are unordered and the distance is Euclidean. The unit
distance problem of Erd\H{o}s \cite{Erdos1946} asks for the order of
growth of $u(n)$. Our main result is the following.

\begin{theorem}\label{thm:main}
There are finite sets $U_j\subset\R^2$, with $|U_j|\to\infty$, such that
\[
 \frac{u(U_j)}{|U_j|^{1.04273}}\longrightarrow\infty.
\]
In particular, $u(n)\ge n^{1.04273}$ for arbitrarily large integers $n$.
\end{theorem}

The largest exponent previously claimed, $1.0418235$, is in the author's
unpublished manuscript \cite{Naslund0418235}. Theorem~\ref{thm:main}
uses the same method with one new ingredient, together with the
corresponding change in the analytic estimate: the class field tower
that supplies the number fields is built over the real quadratic field
$\Q(\sqrt{241})$ instead of over $\Q$. The theorem concerns an unbounded
sequence of cardinalities. Its proof uses no unproved hypothesis, such
as the generalized Riemann hypothesis, but it is computer-assisted:
finite facts about the tower and several numerical inequalities are
certified by exact computation and interval arithmetic, as described in
Section~\ref{intro:code}.

\subsection{Background}\label{intro:background}
Erd\H{o}s \cite{Erdos1946} showed that a suitable section of a scaled
integer lattice gives
\[
 u(n)\ge n^{1+c/\log\log n}
\]
for some $c>0$ and arbitrarily large $n$, and he conjectured that this is
essentially optimal \cite{Erdos1946,Erdos1994}. The best upper bound,
\[
 u(n)=O(n^{4/3}),
\]
is due to Spencer, Szemer\'edi and Trotter
\cite{SpencerSzemerediTrotter1984}; Sz\'ekely \cite{Szekely1997} gave a
proof by crossing numbers.

In $2026$ a team at OpenAI \cite{OpenAI2026} disproved Erd\H{o}s's
conjecture. They constructed sets with $u(U)\ge|U|^{1+\delta}$ for a
fixed $\delta>0$ and arbitrarily large $|U|$, replacing the integer
lattice by ideal lattices in number fields of large degree and small
root discriminant. Alon, Bloom, Gowers, Litt, Sawin, Shankar, Tsimerman,
Wang and Wood \cite{Alon2026} simplified the argument and obtained
$\delta\approx6\cdot10^{-38}$. Sawin \cite{Sawin2026} made every step
explicit and sharpened it, and reached the exponent $1.014114$.

Improvements followed within days in online discussions. On the
Erd\H{o}s Problems forum, mlewko \cite{ErdosProblems90} reached
$1.031849$ in Sawin's criterion with finite data found by ChatGPT, and
on MathOverflow spiderduckpig \cite{SpiderduckpigMO2026} adjusted these
data to reach $1.031883$. In a comment the next day, spiderduckpig
proposed a structural change: in the Golod--Shafarevich step of Sawin's
argument \cite[Lemma~11]{Sawin2026}, replace the unweighted count of
relations by a count weighted by degree in the Zassenhaus filtration.
This was reported to give $\delta\ge0.0333487$
\cite{SpiderduckpigMO2026}. Weighted counts of this kind underlie the
towers of \cite{NaslundLegacy2026}, \cite{Naslund0418235} and the
present paper. The author's MathOverflow answer \cite{NaslundMO2026}
then combined this count with covolume counting, dyadic ramification and
Euler products over a fixed subfield, and reported $\delta\ge0.03583$;
the author's archived manuscript \cite{NaslundLegacy2026} records this
construction with $\delta=0.0358324$. The change log of that answer, of
June~9, 2026, reports a bookkeeping issue affecting the covolume step and
gives $\delta>0.0357$ instead. Table~\ref{intro:table-bounds} lists the
explicit exponents, and the repository \cite{TaoConstants84a} records
further bounds posted online.

\begin{table}[htbp]
\centering
\caption{Explicit lower bounds $u(n)\ge n^{1+\delta}$ for arbitrarily
large $n$, in increasing order. The exponents are rounded down and
constant factors are omitted. The last column gives the kind of source;
apart from Erd\H{o}s's, none of these bounds had appeared in a refereed
journal when this paper was written.}
\label{intro:table-bounds}
\small
\begin{tabular}{@{}llll@{}}
\toprule
Exponent & Source & Main new input & Kind of source\\
\midrule
$1+c/\log\log n$ & Erd\H{o}s \cite{Erdos1946} & integer lattice & journal\\
$1+\delta$, $\delta>0$ & OpenAI \cite{OpenAI2026} & ideal lattices from class field towers & manuscript\\
$1+6\cdot10^{-38}$ & Alon et al.\ \cite{Alon2026} & simplified argument & preprint\\
$1.014114$ & Sawin \cite{Sawin2026} & explicit and sharpened criterion & preprint\\
$1.015263$ & Emmerich \cite{Emmerich2026} & optimized parameters & preprint\\
$1.031185$ & Emmerich, Cordella \cite{EmmerichCordella2026} & extended range of primes & certificate\\
$1.031849$ & mlewko \cite{ErdosProblems90} & finite data found by ChatGPT & forum post\\
$1.0333487$ & spiderduckpig \cite{SpiderduckpigMO2026} & weighted Zassenhaus count & online post\\
$1.0357$ & Naslund \cite{NaslundMO2026,NaslundLegacy2026} & covolume counting, dyadic & online post,\\
 & & ramification, fixed-subfield Euler products & manuscript\\
$1.0418235$ & Naslund \cite{Naslund0418235} & mixed signature, norm-one units, & manuscript\\
 & & nonabelian dyadic local group, profiles & \\
$1.04273$ & this paper & tower over $\Q(\sqrt{241})$ & \\
\bottomrule
\end{tabular}
\end{table}

\subsection{The construction in outline}\label{intro:outline}
The construction carries Erd\H{o}s's lattice argument to number fields
of large degree. Let $K$ be such a field, of degree $2d$, with an
involution $\iota$ whose fixed field is $F$. We take the points of a
translated ideal lattice of $K$ that lie in a bounded region of
$K\otimes\R$. A complex embedding in which $\iota$ acts as complex
conjugation maps them injectively to the plane. It sends every
difference of relative norm one over $F$ to a unit vector. The exponent
is then set by a balance, per unit of $d$. The gain is the number of
equal-norm choices: a prime of $F$ that splits in $K$ and is used with
exponent $k$ gives $k+1$ ideals of the same relative norm. Selecting a
common ideal class loses part of this gain. Through the class-number
formula, that loss is governed by the root discriminant $\lambda$ of $K$
and the value at $s=1$ of $L_F(s)=\zeta_K(s)/\zeta_F(s)$. The geometric
window decides how many of the resulting displacements keep both
endpoints in the region, and its volume fixes the number of points.
Using a prime $\mathfrak p$ with exponent $k$ also makes the ideal
lattice denser by the factor $N\mathfrak p^k$, and the exponent $\delta$
charges the gain for the points this adds.

Quantitatively, suppose that such fields exist with $d\to\infty$, a fixed
root discriminant $\lambda$, fixed local types at the selected primes,
and $\theta=c/d$ bounded below, where $(b,c)$ is the signature of $F$ and
$d=b+2c$; thus $2\theta$ is the proportion of non-real embeddings of $F$.
With hard windows at the selected primes, the exponent $1+\delta$ is
attained when
\begin{equation}\label{intro:margin}
 \sum_q\frac{\log(k_q+1)-\delta k_qf_q\log q}{e_qf_q}
 -\Bigl(\frac12-\delta\Bigr)\log\lambda-C+\mathcal A_\delta(\theta)>0.
\end{equation}
Here $q$ runs over the rational primes below the selected primes, $e_q$
and $f_q$ are their absolute ramification indices and residue degrees in
$F$, $k_q$ is the exponent used at the primes above $q$, $C$ is an upper
bound for $d^{-1}\log L_F(1)$, and
\[
 \mathcal A_\delta(\theta)=(1-\delta)\log2+(1-\theta)\log\pi
 +(1-2\theta)J_\R+\theta J_\C
\]
is the archimedean contribution, in which $J_\R$ and $J_\C$ measure the
overlap of the archimedean profiles at the real and at the complex
places of $F$ under unit displacements (Proposition~\ref{geo:transfer}).
The term $\theta J_\C$ comes from mixed signature: above a complex place
of $F$, the two coordinates of a norm-one element have reciprocal moduli.
Averaging over the norm-one units counts the resulting one-parameter
family of displacements exactly, and the class-number formula cancels the
regulator of these units (Proposition~\ref{geo:mass}).
Section~\ref{fw:section} replaces the hard windows by weighted shells
(Proposition~\ref{fw:transfer}), and Section~\ref{sec:certificate}
verifies the resulting inequality at $\delta=\deltastar$.

The fields must form an infinite family of growing degree with fixed
$\lambda$ and fixed local types. Such families come from infinite
pro-$2$ class field towers with prescribed local Galois groups. The
Golod--Shafarevich criterion \cite{GolodShafarevich} proves that a tower
is infinite when a certain power series in a variable $t$ takes a
negative value; the use of Frobenius conditions to control splitting in
such towers goes back to Hajir, Maire and Ramakrishna
\cite{HajirMaireRamakrishna2021Cutting}.

\subsection{A tower over $\Q(\sqrt{241})$}\label{intro:new}
The fields of \cite{Naslund0418235} come from a pro-$2$ tower over $\Q$
that is unramified outside $2,3,5,7,11,13$, with prescribed local Galois
groups at these primes and Frobenius conditions at the selected primes
up to $31$. Its Golod--Shafarevich series has the form
\[
 1-7t+(\text{a cost for each local condition})-s_D(t).
\]
Here $7$ is the number of generators of the Galois group of the maximal
pro-$2$ extension of $\Q$ unramified outside these six primes, and
$s_D(t)$ is a saving: by Hilbert reciprocity, one dyadic relation is
implied by the others and need not be paid for. Tame ramification at $7$,
$11$ and $13$ supplies generators to pay for the local conditions, and it
costs a factor $\sqrt{1001}$ in the root discriminant.

Now let $B$ be a totally real field of degree $m$ in which every prime
used by the construction splits completely. The generators and the local
conditions of a construction over $\Q$ are then repeated $m$ times over
$B$, while the constant term $1$ and the saving $s_D$ are not; if $P(t)$
is the series over $\Q$, the series over $B$ is roughly
$mP(t)-(m-1)(1-s_D(t))$. The extra room can be spent on ramifying at
fewer primes. The price is the discriminant of $B$, which enters the
root discriminant of the fields. The field $B=\Q(\sqrt{241})$ is the
real quadratic field of smallest discriminant in which $2$, $3$ and $5$
all split, and over it the tower needs no ramification at $7$, $11$ or
$13$. Section~\ref{tw:tower} gives the exact series.

Over $B$ we prescribe, at each of the two primes above $2$, the same
nonabelian dyadic local Galois group of order $32$ as in
\cite{Naslund0418235}; at each of the four primes above $3$ and $5$, a
tame local group $C_2\times C_2$; and conditions that make the Frobenius
elements have order dividing $4$ at the two primes above $29$ and at the
prime $7$, which is inert in $B$. The tower is ramified only above $2$,
$3$ and $5$. Its Galois group has $8$ generators, and the root
discriminant of its fields is
\[
 \lambda=2^{9/4}\sqrt{3615}\approx286.0,
 \qquad 3615=3\cdot5\cdot241,
\]
against $2^{9/4}\sqrt{15015}\approx582.9$ over $\Q$: the base contributes
$\sqrt{241}$, and $\sqrt{1001}$ is saved. The price is paid in the
selected primes: only $2$, $3$, $5$, $7$ and $29$ remain, against the
eleven primes up to $31$ over $\Q$. Theorem~\ref{tw:field-family} states
the resulting family of fields.

The analytic estimate changes accordingly. In \cite{Naslund0418235} the
bound $C$ for $d^{-1}\log L_F(1)$ came from a complete Hecke-factor
calculation in a nonabelian subfield of degree $16384$ of the tower.
Here every field $K$ of Theorem~\ref{tw:field-family} contains the
Kummer field $E=B(\sqrt V)$, of degree $512$ over $\Q$, where $V$ is the
group of $\{2,3,5\}$-units of $B$ modulo squares. Its Dedekind zeta
function is the product of $\zeta_B$ and $255$ quadratic Hecke
$L$-functions of $B$, which we evaluate with certified approximate
functional equations (Sections~\ref{gf:section} and~\ref{an:section}).
Together with refinements from the exact local types of the tower, this
gives $C=\ceiling$ (Theorem~\ref{an:ceiling}).
Table~\ref{intro:table-constructions} compares the two constructions.

\begin{table}[htbp]
\centering
\caption{The construction of \cite{Naslund0418235} and the present one.
The pairs $(e,f)$ are absolute ramification indices and residue degrees
of the selected primes of $F$.}
\label{intro:table-constructions}
\begin{tabular}{lll}
\toprule
 & \cite{Naslund0418235} & this paper\\
\midrule
base field & $\Q$ & $\Q(\sqrt{241})$\\
primes ramified in the tower & $2,3,5,7,11,13$ & above $2,3,5$\\
selected primes, $(e,f)$ & $2{:}\,(8,4)$;\ $3,5{:}\,(2,2)$; & $2{:}\,(8,4)$;\ $3,5{:}\,(2,2)$;\\
 & $7,11,13{:}\,(2,4)$;\ $17,\dots,31{:}\,(1,4)$ & $7{:}\,(1,8)$;\ $29{:}\,(1,4)$\\
root discriminant $\lambda$ & $2^{9/4}\sqrt{15015}\approx582.9$ & $2^{9/4}\sqrt{3615}\approx286.0$\\
$\ell=\log\lambda$ & $6.3680$ & $5.6560$\\
lower bound $\theta_*$ for $\theta$ & $4095/8192$ & $65535/131072$\\
ceiling $C$ & $0.042161819$ & $0.04871285$\\
exponent $1+\delta$ & $1.0418235$ & $1.04273$\\
certified margin & $5.3\cdot10^{-6}$ & $1.77\cdot10^{-4}$\\
\bottomrule
\end{tabular}
\end{table}

\subsection{Why the exponent improves}\label{intro:why}
In \eqref{intro:margin}, lowering $\log\lambda$ from $6.3680$ to
$5.6560$ is worth $(1/2-\delta)\cdot0.712\approx0.33$ per unit of $d$ at
$\delta\approx0.043$. Three terms pay for it. The first sum, the gain
from equal-norm choices net of their cost in vertices, falls from $1.034$
to $0.721$, because the tower over $B$ selects only five rational primes,
and two of them, $7$ and $29$, with large residue degrees. The ceiling
$C$ rises by $0.0066$, because the Kummer field carries less Frobenius
information than the field of degree $16384$ used over $\Q$. The
archimedean term falls by $0.013$, mostly because $\delta$ is larger.
Table~\ref{intro:table-budget} shows the certified terms of the margin
for both constructions, each at its own exponent.

\begin{table}[htbp]
\centering
\caption{The terms of the certified margin, per unit of $d$, rounded to
seven decimals. The first row is the weighted-shell form of the first sum
in \eqref{intro:margin} (Section~\ref{fw:section}), and the archimedean
term is computed with the certified lower bound $J_\C^*$ for $J_\C$. The
margin is the quantity $M_*(\theta_*)$ of Section~\ref{sec:certificate}
and of the certificate of \cite{Naslund0418235}.}
\label{intro:table-budget}
\begin{tabular}{lrr}
\toprule
Term & \cite{Naslund0418235}, $\delta=0.0418235$ & this paper, $\delta=\deltastar$\\
\midrule
finite sum & $1.0335669$ & $0.7214955$\\
$-(1/2-\delta)\log\lambda$ & $-2.9176605$ & $-2.5863213$\\
$-C$ & $-0.0421618$ & $-0.0487129$\\
$\mathcal A_\delta(\theta_*)$ & $1.9262607$ & $1.9137154$\\
\midrule
margin & $0.0000053$ & $0.0001767$\\
\bottomrule
\end{tabular}
\end{table}

The balance leaves a certified margin of about $1.77\cdot10^{-4}$ per unit
of $d$ at $\delta=\deltastar$. With the same profiles, and shell weights
optimized in the same way, the margin at $\delta=0.0428$ is negative.

\subsection{Relation to Sawin's criterion}\label{intro:comparison}
All bounds in Table~\ref{intro:table-bounds} from Sawin's onward use his
framework \cite{Sawin2026}: a quadratic extension $K/F$, ideals of equal
relative norm built from split primes, a pigeonhole over ideal classes,
and fields from Golod--Shafarevich towers with Frobenius conditions.
Section~1.6 of \cite{Naslund0418235} compares that criterion with
\eqref{intro:margin} term by term. In brief, counting vertices by
covolume instead of packing returns $\delta\log\lambda$. The exact
class-number formula, with an explicit Euler-product bound $C$, replaces
Louboutin's bound for the relative class number \cite{Louboutin2000}
together with the factor $2^d$ that Sawin allows for units; their
combined cost $1+\log\log\lambda$ would be about $2.73$ here. Mixed
signature contributes the term $\theta(J_\C-2J_\R-\log\pi)$, worth more
than $0.31$ per unit of $d$ here. Optimized profiles and shells replace
the sup-norm window. The present paper keeps all of these and changes the
tower. Sawin suggested that towers satisfying different local conditions
might improve the exponent \cite{Sawin2026}. The base field is a change
of a different kind: it adds no local conditions at new primes, but
repeats the conditions at the small primes over two primes of $B$ each.

\subsection{Organization}\label{intro:organization}
Section~\ref{tw:tower} constructs the tower over $B$ and proves
Theorem~\ref{tw:field-family}, which supplies the fields.
Section~\ref{gf:section} describes the Kummer field and its quadratic
$L$-functions, and Section~\ref{an:section} proves the ceiling
$d^{-1}\log L_F(1)<C$ (Theorem~\ref{an:ceiling}).
Sections~\ref{geo:section}--\ref{pf:section} carry out the geometric
construction: the norm-one mass formula and the transfer to planar sets,
the weighted windows at the selected primes, and the archimedean
profiles. These sections hold for any quadratic extension with the
stated properties and follow \cite{Naslund0418235} closely.
Section~\ref{sec:certificate} evaluates the final inequality at
$\delta=\deltastar$ and completes the proof of Theorem~\ref{thm:main}.

\subsection{Computations and data}\label{intro:code}
The proof uses finite computations: linear algebra over $\Ftwo$ for the
tower, exact rational arithmetic for the Golod--Shafarevich series,
certified evaluations of $L$-functions in the ball arithmetic of Arb,
now part of FLINT \cite{Johansson2017,FLINT}, and directed interval
bounds for the geometric margin \cite{Mpmath}. PARI/GP \cite{PARI} is
used for the arithmetic of number fields and for independent checks.
Each finite fact is verified by a supplementary program, named in the
text where the fact is used. These certificates concern exact arithmetic
and rigorous inequalities, not the statistical accuracy of
floating-point approximations.

The programs and their data are in the directory
\texttt{papers/0.04273/certificates} of the public repository
\cite{NaslundCode2026}. Some certified routines, for the degree-two
approximate functional equations, the archimedean profiles and the
finite windows, are imported unchanged from the supplementary archive of
\cite{Naslund0418235}, which is in the same repository together with a
manifest of the SHA-256 hashes of its members. A replay program runs
every step and records its results. The computations were run with
PARI/GP~2.17.2, python-flint~0.9.0 (FLINT~3.6.0) and mpmath~1.3.0.

A separate Lean~4 development \cite{NaslundLean2026}, built on Mathlib
\cite{Lean4_2021,Mathlib2020}, proves the planar conclusion with the
slightly smaller exponent $1.0427$ from one explicit numerical inequality
for the Dedekind zeta function of $E$
(Remark~\ref{an:lean-hypothesis}). That inequality has not been proved in
Lean, so the formalization is conditional, and it is not used as a
substitute for the arguments here.
